\documentclass[11pt]{article}
\usepackage{mathblatt}
\usepackage{multicol}


\begin{document}
\pfheftstil
\fancyfoot[C]{\parbox[t]{\textwidth}{\mbfhbox\par\vspace{1.5mm}{\footnotesize\color{mbgrau}{\scriptsize MRD}\hfill\thepage}}}
\pfheftkopf{Satz des Pythagoras \textperiodcentered{} Lösungen}{}
{\small\noindent\textbf{Typische Fehler}\par\smallskip\noindent \textbullet\ Quadrate addiert, obwohl eine Kathete gesucht ist: BC² = 9² + 4² (9,85 m statt 8,06 m); 384² + 255² (461 m statt 287,1 m); s² + r² für die Kegelhöhe.\par\noindent \textbullet\ Falsche Seite als Hypotenuse: die längere Kathete für die Hypotenuse gehalten (√(34² $-$ 12²) = 31,8 m statt 36,1 m); die Stufenhöhe der Rampe als Hypotenuse; x² = y² + z² angekreuzt.\par\noindent \textbullet\ Wurzel vergessen: z = y² + x² oder z = x² $-$ y² angekreuzt; c² als Ergebnis angegeben.\par\noindent \textbullet\ Satz als Längenformel c = a + b; „Pythagoras heißt Summe“ gemerkt und darum bei der Kathete die Differenz nicht erkannt.\par\noindent \textbullet\ Wurzel aus jedem Summanden einzeln: √(a² + b²) = a + b.\par\noindent \textbullet\ Ganze statt halbe Seite: a statt a/2 im Stützdreieck der Pyramide; die halbe Diagonale mit der ganzen (AF = 50 statt 25) angesetzt; 30 statt 15 als Überstand am Trapez.\par\noindent \textbullet\ Radius und Durchmesser vertauscht: Radius als Kathete im Becher ($\approx$ 9,7 statt 11,5 cm); Durchmesser als Kathete im Kegeldach.\par\noindent \textbullet\ Schräge als Höhe: Mantellinie als Kegelhöhe (Gesamthöhe 33,6 statt 32,2 m); Trapezhöhe als Schenkel.\par\noindent \textbullet\ Buchstabenfixierung: „c ist immer die Hypotenuse“ – bei Beschriftung mit x, y, z oder wenn c eine Kathete ist, wird die Formel unverändert übernommen.\par}\medskip
\begin{pfloesung}{}
\lz{1.}{$c^2 = a^2 + b^2$}{}{}
\end{pfloesung}
\begin{pfloesung}{}
\lz{2.}{a) $b^2 = a^2 + c^2$; b) $x^2 = y^2 + z^2$; c) $q^2 = p^2 + s^2$; d) $v^2 = u^2 + w^2$}{}{}
\end{pfloesung}
\begin{pfloesung}{}
\lz{3.}{z² = x² + y²}{}{}
\end{pfloesung}
\begin{pfloesung}{}
\lz{4.}{Die Summe der Flächeninhalte der Kathetenquadrate ist gleich dem Flächeninhalt des Hypotenusenquadrates}{a² + b² = c²}{}
\end{pfloesung}
\begin{pfloesung}{}
\lz{5.}{b) Hypotenuse, $s^2 = 3^2 + 0{,}4^2$; c) Hypotenuse, $d^2 = 48^2 + 27^2$; d) Kathete, $25^2 = h^2 + 7^2$}{}{}
\end{pfloesung}
\begin{pfloesung}{}
\lz{6.}{z = √(x² + y²)}{}{}
\end{pfloesung}
\begin{pfloesung}{}
\lz{7.}{$c = 39$ cm}{Gleichung: $c^2 = 36^2 + 15^2$; Zwischenergebnis: $c^2 = 1\,296 + 225 = 1\,521$; Wurzel: $c = \sqrt{1\,521} = 39$ cm}{}
\end{pfloesung}
\begin{pfloesung}{}
\lz{8.}{$\overline{AB}$ = 10√13 $\approx$ 36,1 m}{$\overline{AB}^2$ = 12² + 34² = 1300}{}
\end{pfloesung}
\begin{pfloesung}{}
\lz{9.}{h = 4,8 m}{h² = 5² $-$ 1,4² = 23,04}{}
\end{pfloesung}
\begin{pfloesung}{}
\lz{10.}{$\overline{FA}$ $\approx$ 287,1 m}{$\overline{FA}$² = 384² $-$ 255² = 82 431}{}
\end{pfloesung}
\begin{pfloesung}{}
\lz{11.}{$\overline{AD}$ = √553 719 $\approx$ 744 m}{$\overline{AD}^2$ = 920² $-$ 541²}{}
\end{pfloesung}
\begin{pfloesung}{}
\lz{12.}{Katheten: die Höhe und der Überstand; Hypotenuse: der Schenkel}{}{}
\end{pfloesung}
\begin{pfloesung}{}
\lz{13.}{$h = 63$ cm}{halbe Grundseite: $\tfrac{32}{2} = 16$ cm; Gleichung umstellen: $h^2 = 65^2 - 16^2$; Zwischenergebnis: $h^2 = 4\,225 - 256 = 3\,969$; Wurzel: $h = \sqrt{3\,969} = 63$ cm}{}
\end{pfloesung}
\begin{pfloesung}{}
\lz{14.}{$\overline{BC}$ = √20 $\approx$ 4,47}{$\overline{AB}$ = 2, $\overline{AC}$ = 4; $\overline{BC}$² = 2² + 4² = 20}{}
\end{pfloesung}
\begin{pfloesung}{}
\lz{15.}{$\overline{AB}$ $\approx$ 60,4 cm}{$\overline{AF}$ = 25 cm, $\overline{BF}$ = 80 $-$ 25 = 55 cm; $\overline{AB}$² = 25² + 55² = 3650}{}
\end{pfloesung}
\begin{pfloesung}{}
\lz{16.}{Gesamthöhe $\approx$ 32,2 m}{Kegelhöhe² = 8,6² $-$ 4,7² = 51,87; Kegelhöhe $\approx$ 7,20 m}{}
\end{pfloesung}
\begin{pfloesung}{}
\lz{17.}{√89,96 + 2 $\approx$ 11,5 cm}{Diagonale √(7² + 6,4²) = √89,96 $\approx$ 9,5 cm; Durchmesser 6,4 cm}{}
\end{pfloesung}
\end{document}